\section{File formats, canonical form and test vectors}\label{sec:formats}


This appendix collects what an independent implementer needs in order to read the two certificate files
and to test a verifier written from \cref{alg:verify,alg:leaf}. Cells are named by their \emph{ids}: the
cell $(x,y,z)\in\{0,\dots,5\}^3$ has the id $36x+6y+z\in\{0,\dots,215\}$ (\cref{ssec:canon}), and all files
use ids. The lines and the group $\G$ are those of \cref{sec:board}; a summary of the geometry in terms of
ids, with the ids of the $24$ busy cells, is given in \cref{ssec:fmt-cells}. Both certificate files are
plain ASCII text files. Every line,
including the last one, is terminated by a single line feed (byte \texttt{0x0A}); there are no carriage
returns, no empty lines, no tabs and no spaces other than the single spaces that separate the fields;
the only characters that occur are the digits, the comma, the hyphen, the space and the line feed.
Numbers are written in decimal without leading zeros.

\subsection{Sets, their order, and canonical forms}\label{ssec:fmt-canon}

A set of cells is written as the list of its ids in ascending order, separated by commas and without
spaces; the empty set is written as a single hyphen. Sets of the same size are ordered
lexicographically by these lists: $A<A'$ if, at the first position where the two ascending lists differ,
the id in the list of $A$ is the smaller one. (Equivalently: the smallest element of the symmetric
difference of $A$ and $A'$ belongs to $A$.)

The \emph{canonical form} of a pair $(B,F)$ of disjoint sets is the pair $(g(B),g(F))$, $g\in\G$, that is
smallest when the images of $B$ are compared first and, among the maps that minimize $g(B)$, the images
of $F$ second. A state $\Pst(B,F,t)$ or $\Qst(B,F,t)$ is canonical if $(B,F)$ is; $t$ and the type of the
state are not affected. Both files contain canonical states only. If $(B,F)$ is any state and $g$ is a map
with $(g(B),g(F))$ canonical, the table move $m$ listed for the canonical Q-state corresponds to the move
$g^{-1}(m)$ in the frame of $(B,F)$; any $g$ attaining the canonical form may be used (\cref{lem:canon}).

For a verifier it suffices to compute the canonical form by trying all $1{,}152$ maps. Test vectors (ids;
each line reads $B\mid F\ \to\ $ canonical $B\mid$ canonical $F$):
\begin{center}\small\ttfamily
\begin{tabular}{@{}l@{}}
43 | 0 \ $\to$\ 0 | 43\\
0,46,92,123,169,180,215 | - \ $\to$\ 0,5,61,93,122,154,210 | -\\
0,35,46,92,123,169,180,215 | - \ $\to$\ 0,5,61,93,122,154,210,215 | -\\
5,30,43,64,86,93,128,169,180 | 0,35,46,87,122,151,172,185,210\\
\qquad$\to$\ 0,35,43,64,86,93,123,154,210 | 5,30,61,122,129,151,172,180,215
\end{tabular}
\end{center}
The second and third are the sets $\Tseven$ and $\Teight$ of \cref{ex:T7,ex:T8}; the last is the final
position of Play~B in \cref{ssec:walk}.

\subsection{The strategy table}\label{ssec:fmt-table}

The file \fn{results/final/strategy.tbl} is a text file with $93{,}164$ lines (its \emph{entries}), one
for each Q-state of the
certificate, in no particular order and without a header. An entry has four fields separated by single
spaces:
\begin{center}
\texttt{t\ \ B\ \ F\ \ m}
\end{center}
where $t$ is the round in decimal, $B$ and $F$ are sets in the notation above, and $m$ is a cell id. The
entry says: \emph{in the canonical Q-state $\Qst(B,F,t)$, in which Maker has made $t$ moves, Breaker owns the
$t-1$ stones $B$, and the stones $F$ of Maker are revealed ($|F|\le t$), the table move of Breaker is the
cell $m$}.
The move is given in the frame of the canonical state and satisfies $m\notin B\cup F$. No key $(t,B,F)$
occurs twice. Three entries of the file (the first, the $2{,}354$th and the $9{,}791$st):
\begin{verbatim}
1 - - 0
5 0,46,169,215 - 92
7 0,35,46,92,123,169 5,30,43,180,185,210,215 50
\end{verbatim}
In coordinates: the first is Breaker's opening move $(0,0,0)$ in the state $\Qst(\emptyset,\emptyset,1)$.
The second says that with $B=\{(0,0,0),(1,1,4),(4,4,1),(5,5,5)\}$ and nothing revealed, Breaker plays
$(2,3,2)$ in round $5$. In the third, Breaker owns $(0,0,0)$, $(0,5,5)$, $(1,1,4)$, $(2,3,2)$, $(3,2,3)$,
$(4,4,1)$, the seven revealed stones of Maker are $(0,0,5)$, $(0,5,0)$, $(1,1,1)$, $(5,0,0)$, $(5,0,5)$,
$(5,5,0)$, $(5,5,5)$, and Breaker's move in round $7$ is the non-busy cell $(1,2,2)$.

\subsection{The leaf file}\label{ssec:fmt-leaves}

The file \fn{results/final/leaves.txt} has $106{,}723$ lines (entries) with three fields,
\begin{center}
\texttt{t\ \ B\ \ F}
\end{center}
meaning: \emph{the canonical P-state $\Pst(B,F,t)$, in which $t$ rounds are complete, Breaker owns the $t$
stones $B$, and the stones $F$ of Maker are revealed, is a leaf}, i.e.\ $\LEAF(B,F,t)$ is claimed to hold
and the
state is not expanded. Three entries (the first, the $8{,}185$th and the last):
\begin{verbatim}
9 0,1,5,30,43,93,122,172,210 7,13,14,19,21,25,28,31,35
8 0,5,61,93,122,154,210,215 -
9 0,35,43,64,86,93,123,154,210 5,30,61,122,129,151,172,180,215
\end{verbatim}
In the first, Breaker owns $(0,0,0)$, $(0,0,1)$, $(0,0,5)$, $(0,5,0)$, $(1,1,1)$, $(2,3,3)$, $(3,2,2)$,
$(4,4,4)$, $(5,5,0)$, and $F$ consists of the nine cells $(0,1,1)$, $(0,2,1)$, $(0,2,2)$, $(0,3,1)$,
$(0,3,3)$, $(0,4,1)$, $(0,4,4)$, $(0,5,1)$, $(0,5,5)$: Maker has made two threats, on the line $\{(0,i,1)\}$ and on the diagonal
$\{(0,i,i)\}$ of the face $x=0$, which share the stone $(0,1,1)$; Breaker's forced replies are $(0,0,1)$ and
$(0,0,0)$, and all nine stones of Maker are revealed ($\theta=0$). The second entry is the canonical form of
$\Pst(\Teight,\emptyset,8)$, the only leaf with $\theta=8$. The third is the final state of Play~B of
\cref{ssec:walk}.

A verifier must not assume that the files are correct; \cref{alg:verify} only uses them as hints. An
entry is \emph{well formed} if it has the right number of fields, $t$ is a non-negative integer, $B$ and
$F$ are strictly ascending lists of ids between $0$ and $215$ (or the hyphen), $B\cap F=\emptyset$,
$|F|\le t$, $|B|=t-1$ for a Q-state and $|B|=t$ for a P-state, and, in the table, $m$ is an id with
$m\notin B\cup F$. What a verifier has to do with an entry that is not well formed, with an entry whose
state is not canonical, with a key that occurs twice and with an entry that is never reached is stated in
\cref{ssec:what}: none of these defects can make a sound verifier accept wrongly, but a verifier that is
to establish \cref{fact:H} as stated must reject or report all four. Both files consist of well-formed,
canonical, pairwise different entries, all of which are reached.

\subsection{Numbers a verifier should reproduce}\label{ssec:fmt-targets}

A correct implementation of \cref{alg:verify} with $\Rmax=10$ accepts the two files and visits
$246{,}299$ states: $153{,}135$ P-states, of which $46{,}412$ are expanded and $106{,}723$ are leaves, and
$93{,}164$ Q-states. The numbers of P-states, of leaves, of Q-states, of threat children, of danger
children and of \VOID{} children of type (ii) in each round $t=0,\dots,10$ are those of \cref{tab:tree};
they are printed once more, with one row for each kind of number, in \cref{ssec:supp-targets}.
The conventions are those of \cref{tab:tree} and of the comment ``Counting'' in \cref{ssec:algverify}:
states are counted as canonical forms; a threat or danger
child is counted, at the round of the child, once for every expanded canonical P-state of which it is a
non-\VOID{} child. By \cref{lem:children-distinct} different data $(l,e)$, respectively
$(\{a,b\},Q_a,Q_b)$, of one state define different children, so this is the number of data for which
\cref{alg:verify} generates a non-\VOID{} child; children of one state may be symmetric to each other or
to children of other states, and they are then counted separately. A \VOID{} child of type (ii) is
counted at the round of its parent Q-state. In total there are $143{,}498$ threat children, $28{,}275$ danger children
and $54{,}482$ children of type (ii), of which $151$ are \VOID; no child of type (c) or (d) is \VOID. Every
entry of both files is reached. The leaves split by $\theta=t-|F|=0,\dots,8$ into
$64{,}841$, $32{,}703$, $7{,}267$, $1{,}434$, $343$, $105$, $25$, $4$, $1$ (\cref{tab:leaf-theta}), and $434$ table moves are
non-busy cells ($16$, $399$, $11$, $8$ in rounds $7$ to $10$).

\subsection{Test vectors for the leaf test}\label{ssec:fmt-leaf}

Let
\[
\Tseven=\{(0,0,0),(5,5,5),(1,1,4),(4,4,1),(2,3,2),(3,2,3),(5,0,0)\},\qquad
\Teight=\Tseven\cup\{(0,5,5)\},
\]
with ids \texttt{0,46,92,123,169,180,215} and \texttt{0,35,46,92,123,169,180,215}. \Cref{tab:leafvectors}
lists values of $\LEAF(B,F,k)$ that an implementation of \cref{alg:leaf} should reproduce; in the
notation of the files an instance is the line ``$k$ $B$ $F$'' formed from the first three columns. Each
value was obtained with three independent
implementations (the two verifiers and the library of the search), which also agree on which of the
three conditions (L1)--(L3) of \cref{def:leaf} fails. During the internal refereeing of the first version
of this paper (\cref{ssec:impl}) all $26$ instances were
reproduced, including the failing condition, by three further implementations of \cref{alg:leaf}. The last column gives, for comparison, the number
$\worst(B,F,k)$ of \cref{sec:ip} computed with the integer program; the pairing condition
$\CC(B,F,k)$ holds if and only if it is $0$.

\begin{table}[ht]
\centering\footnotesize\setlength{\tabcolsep}{3pt}
\begin{tabular}{@{}r>{\raggedright\arraybackslash}p{0.265\textwidth}>{\raggedright\arraybackslash}p{0.2\textwidth}>{\raggedright\arraybackslash}p{0.275\textwidth}lr@{}}
\toprule
$k$ & $B$ (ids) & $F$ (ids) & description & $\LEAF$ & $\worst$\\
\midrule
\texttt{7} & \fn{0,46,92,123,169,180,215} & \texttt{-} & $\Tseven$, $k=7$ & false (3) & 6\\
\texttt{8} & \fn{0,35,46,92,123,169,180,215} & \texttt{-} & $\Teight$, $k=8$ & true & 0\\
\texttt{9} & \fn{0,35,46,92,123,169,180,215} & \texttt{-} & $\Teight$, $k=9$ & false (3) & 1\\
\texttt{8} & \fn{0,35,46,92,123,169,180,215} & $h$ & $\Teight$, $F=\{h\}$; one instance for each of the
  $16$ busy cells $h\notin\Teight$ & true & 0\\
\texttt{8} & \fn{0,35,43,46,92,123,169,180} & \texttt{215} & $\Teight-(5,5,5)+(1,1,1)$, $F=\{(5,5,5)\}$ & true & 0\\
\texttt{8} & \fn{0,5,35,43,46,61,86,92} & \fn{123,129,169,172,180,210,215} & antipode stealing after $8$
  rounds & false (2) & 6\\
\texttt{9} & \fn{0,5,35,43,46,61,86,87,92} & \fn{123,129,154,169,172,180,210,215} & antipode stealing after
  $9$ rounds & true & 0\\
\texttt{10} & \fn{0,5,30,35,43,46,61,86,87,92} & \fn{123,128,129,154,169,172,180,210,215} & antipode
  stealing after $10$ rounds & true & 0\\
\texttt{9} & \fn{0,8,35,46,92,123,169,180,215} & \fn{44,80,116,152,188} & threat: $\Teight\cup\{(0,1,2)\}$,
  $F=\{(x,1,2)\colon x\ge1\}$ & true & 0\\
\texttt{8} & \fn{0,8,46,92,123,169,180,215} & \fn{44,80,116,152,188} & threat: $\Tseven\cup\{(0,1,2)\}$, the
  same $F$ & false (3) & 2\\
\texttt{0} & \texttt{-} & \texttt{-} & the empty board & false (2) & 80\\
\bottomrule
\end{tabular}
\caption{Test vectors for $\LEAF(B,F,k)$, with $B$ and $F$ in separate columns. ``false ($i$)'' means that
condition (L$i$) of \cref{def:leaf}
is the first to fail: (L2) no pairing for $(B,F)$; (L3) an $F$-connected set of at least three live lines
has surplus less than $k-|F|$. The fourth row stands for $16$ instances; the $16$ cells $h$ have the ids
5, 30, 43, 61, 64, 86, 87, 93, 122, 128, 129, 151, 154, 172, 185, 210. In all there are $26$ instances.}
\label{tab:leafvectors}
\end{table}

Some explanations (see also \cref{fact:T8-leaf}). For $\Tseven$ the set of all $104$ live lines is
$\emptyset$-connected and has surplus $216-7-208=1=\slack(\Tseven)<7$. For $\Teight$ the $100$ live lines
have surplus $216-8-200=8$, which is the minimum over all $\emptyset$-connected sets of at least three
live lines; hence $\LEAF(\Teight,\emptyset,8)$ holds and $\LEAF(\Teight,\emptyset,9)$ fails. In the ``antipode stealing''
instances Breaker has played $(0,0,0)$, $(1,1,4)$, $(2,3,2)$, $(0,5,5)$, $(1,1,1)$, $(2,2,2)$, $(0,0,5)$,
$(1,4,1)$ and then $(2,2,3)$ and $(0,5,0)$, and $F$ consists of the antipodes of all these stones except the
last one. In the two threat instances Maker holds five cells of the line $\{(x,1,2)\}$ and Breaker has been
forced to $(0,1,2)$. On the empty board $148$ lines would need $296$ cells, and the deficiency
$296-216=80$ is the value of $\worst$.

Further values can be generated with the programs of the archive: \fn{bin/phicheck} and
\fn{bin/treecheck} with the option \fn{--leafcheck} read instances ``$k$ $B$ $F$'' from standard input and
print the value of $\LEAF$, the failing condition and a witness.

\subsection{Test vectors for \texorpdfstring{$\LOST$}{LOST}}\label{ssec:fmt-lost}

The predicate $\LOST$ of \cref{def:lost} is false at every expanded state of the certificate
(\cref{prop:lost-stones}), so that an implementation of it cannot be tested on the certificate; a verifier
for $\Rmax=10$ does not need it. For implementers who want it nevertheless, for instance for experiments
with larger $\Rmax$, \cref{tab:lostvectors} gives nine instances $(B,F,t)$ with $t\ge 9$. The cell $x$
of a pattern may belong to $F$ (see the comment after \cref{def:lost}). The values were computed with
\fn{bin/treecheck} \fn{--lostcheck}, which reads lines ``$t$ $B$ $F$'' from standard input, and
cross-checked with two other implementations: the function \texttt{lost} of \fn{verify/spotcheck.py} and
the literal test of \cref{def:lost} in the verifier that was written from the paper alone during the
internal refereeing
(\fn{logs/paper/replication_from_paper/verify.cpp}). The three agree on the verdict and on which kind of
pattern fits; the instances, the outputs and a script that repeats the comparison are in
\fn{logs/paper/lost_vectors/}. The first and the last implementation agree, moreover, on the $800$
instances with $t=9,\dots,13$ that were used in the audit of the main verifier (\cref{ssec:supp-audits}).

\begin{table}[ht]
\centering\footnotesize\setlength{\tabcolsep}{3pt}
\begin{tabular}{@{}r>{\raggedright\arraybackslash}p{0.33\textwidth}>{\raggedright\arraybackslash}p{0.33\textwidth}ll@{}}
\toprule
$t$ & $B$ (ids) & $F$ (ids) & $\LOST$ & fitting patterns\\
\midrule
\texttt{9} & \texttt{-} & \texttt{-} & false & none\\
\texttt{10} & \texttt{-} & \texttt{-} & true & $\Etwo$ ($11$ stones)\\
\texttt{10} & \fn{0,56,61,65,91,92,146,151,154,172} & \texttt{-} & true & $\Etwo$\\
\texttt{10} & \fn{5,30,98,123,128,129,141,169,172,185} & \fn{15,79,148} & false & none\\
\texttt{10} & \fn{1,21,50,128,135,136,180,183,206,210} & \fn{5,23,59,77,83,107,131,173,203,215} & true & $\Etwo$\\
\texttt{11} & \fn{43,46,61,64,87,149,169,185,209,210,215} & \fn{7,83,204} & false & none\\
\texttt{12} & \fn{12,23,24,30,87,88,142,156,165,172,180,210} & \fn{40,43,45,46,48,50,53,58,63,64,69,71} & true
  & $\Ethree$ only ($13$ stones)\\
\texttt{12} & \fn{5,11,79,86,92,93,94,121,123,154,186,210} & \fn{26,70,77,194} & false & none\\
\texttt{10} & \fn{54,55,127,131,133,142,143,157,164,214} & \fn{0,12,18,90,108,120,144,150,180} & true
  & $\Etwo$ ($x\in F$)\\
\bottomrule
\end{tabular}
\caption{Test vectors for $\LOST(B,F,t)$. In the notation of the files an instance is the line
``$t$ $B$ $F$''. The last column says which kinds of patterns fit $(B,F,t)$; ``$\Ethree$ only'' means that no
$\Etwo$ pattern fits. $\LOST(B,F,t)$ is defined for arbitrary disjoint sets $B$, $F$ and $t\ge 0$
(\cref{def:lost}); in the first two instances $t$, $B$ and $F$ are such arbitrary inputs of the predicate and
not states ($|B|\ne t$), and they must be read without the test $|B|=t$ of \cref{ssec:fmt-leaves}. In the
other seven instances $|B|=t$ and $F$ is $B$-admissible. In the fifth and in the seventh instance all stones
of Maker are revealed. In the last instance every fitting pattern has $x\in F$; an implementation that
requires $x\notin F$ returns false there.}
\label{tab:lostvectors}
\end{table}
